% GENERATED FILE. Edit canonical lessons, then run npm run generate:books.
% canonical-source-hash: fnv1a64:f22e8f6641339087
% canonical-lessons: AR-C481-zakhkha, AR-C481-radhadh, AR-C481-wabil, AR-C481-hatala, AR-C481-safa

\chapter{Rain}
\label{ch:ar-rain}
\hlchaptermodality{481}

\begin{chapteropening}
\textbf{By the end of this chapter:} \emph{I can say a shower of rain, drizzle, a downpour, to pour down and to clear up.}

Complete the last lesson of chapter 481: I can say a shower of rain, drizzle, a downpour, to pour down and to clear up.
\end{chapteropening}

\section[\texorpdfstring{\ar{زخة} (zakhkha)}{زخة (zakhkha)}]{\ar{زخة} (zakhkha) — a shower of rain}

\label{lesson:AR-C481-zakhkha}



\begin{quote}
Before the new one: say the Arabic for to drop, then the Arabic for mild.
\end{quote}

\subsection*{You'll want to know: \ar{زخة}}
\textbf{\ar{زخة}} — \emph{zakhkha} — \textquotedblleft{}a shower of rain\textquotedblright{}.

\begin{quote}
\ar{نزلت} \ar{زخة} \ar{مطر} \ar{قصيرة}.
\end{quote}

\begin{quote}
\emph{nazalat zakhkhat maṭar qaṣīra.}
\end{quote}

\begin{quote}
A short shower of rain fell.
\end{quote}

\begin{cousinweb}
\textbf{\ar{زخة}} (\emph{zakhkha}) names a short, sudden shower; the origin is not settled.
\end{cousinweb}

\begin{grammarlens}[title={What you've built}]
One more word for getting around.
\end{grammarlens}

\subsection*{Your turn}
\begin{itemize}
\raggedright
  \item \emph{Say it:} \emph{zakhkha}
  \item \emph{Say it:} \emph{zakhkha}, once more
  \item \emph{Say it:} \emph{nazalat zakhkhat maṭar qaṣīra.} — A short shower of rain fell.
  \item \emph{Recall:} say the Arabic for to drop, then the Arabic for mild, then say \emph{zakhkha} again
\end{itemize}

\begin{tcolorbox}[breakable,colback=teal!4,colframe=teal!35!black,title={Before you move on}]
What does \ar{زخة} mean? (\textquotedblleft{}A shower of rain\textquotedblright{}.) Say it once more.
\end{tcolorbox}
\section[\texorpdfstring{\ar{رذاذ} (radhādh)}{رذاذ (radhādh)}]{\ar{رذاذ} (radhādh) — drizzle; spray}

\label{lesson:AR-C481-radhadh}



\begin{quote}
Before the new one: say the Arabic for mild, then the Arabic for a shower of rain.
\end{quote}

\subsection*{You'll want to know: \ar{رذاذ}}
\textbf{\ar{رذاذ}} — \emph{radhādh} — \textquotedblleft{}drizzle; spray\textquotedblright{}.

\begin{quote}
\ar{كان} \ar{الرذاذ} \ar{خفيفا} \ar{في} \ar{الصباح}.
\end{quote}

\begin{quote}
\emph{kāna ar-radhādh khafīfan fī aṣ-ṣabāḥ.}
\end{quote}

\begin{quote}
The drizzle was light in the morning.
\end{quote}

\begin{cousinweb}
\textbf{\ar{رذاذ}} (\emph{radhādh}) is usually traced to the native Arabic root \emph{r-dh-dh}, \textquotedblleft{}fine rain\textquotedblright{}.
\end{cousinweb}

\begin{grammarlens}[title={What you've built}]
One more word for getting around.
\end{grammarlens}

\subsection*{Your turn}
\begin{itemize}
\raggedright
  \item \emph{Say it:} \emph{radhādh}
  \item \emph{Say it:} \emph{radhādh}, once more
  \item \emph{Say it:} \emph{kāna ar-radhādh khafīfan fī aṣ-ṣabāḥ.} — The drizzle was light in the morning.
  \item \emph{Recall:} say the Arabic for mild, then the Arabic for a shower of rain, then say \emph{radhādh} again
\end{itemize}

\begin{tcolorbox}[breakable,colback=teal!4,colframe=teal!35!black,title={Before you move on}]
What does \ar{رذاذ} mean? (\textquotedblleft{}Drizzle; spray\textquotedblright{}.) Say it once more.
\end{tcolorbox}
\section[\texorpdfstring{\ar{وابل} (wābil)}{وابل (wābil)}]{\ar{وابل} (wābil) — a downpour}

\label{lesson:AR-C481-wabil}



\begin{quote}
Before the new one: say the Arabic for a shower of rain, then the Arabic for drizzle.
\end{quote}

\subsection*{You'll want to know: \ar{وابل}}
\textbf{\ar{وابل}} — \emph{wābil} — \textquotedblleft{}a downpour\textquotedblright{}.

\begin{quote}
\ar{نزل} \ar{وابل} \ar{من} \ar{المطر} \ar{في} \ar{المساء}.
\end{quote}

\begin{quote}
\emph{nazala wābil min al-maṭar fī al-masāʾ.}
\end{quote}

\begin{quote}
A downpour came in the evening.
\end{quote}

\begin{cousinweb}
\textbf{\ar{وابل}} (\emph{wābil}) is built on the native Arabic root \emph{w-b-l}, \textquotedblleft{}heavy rain of large drops\textquotedblright{} (Lane).
\end{cousinweb}

\begin{grammarlens}[title={What you've built}]
One more word for getting around.
\end{grammarlens}

\subsection*{Your turn}
\begin{itemize}
\raggedright
  \item \emph{Say it:} \emph{wābil}
  \item \emph{Say it:} \emph{wābil}, once more
  \item \emph{Say it:} \emph{nazala wābil min al-maṭar fī al-masāʾ.} — A downpour came in the evening.
  \item \emph{Recall:} say the Arabic for a shower of rain, then the Arabic for drizzle, then say \emph{wābil} again
\end{itemize}

\begin{tcolorbox}[breakable,colback=teal!4,colframe=teal!35!black,title={Before you move on}]
What does \ar{وابل} mean? (\textquotedblleft{}A downpour\textquotedblright{}.) Say it once more.
\end{tcolorbox}
\section[\texorpdfstring{\ar{هطل} (haṭala)}{هطل (haṭala)}]{\ar{هطل} (haṭala) — to pour down (of rain)}

\label{lesson:AR-C481-hatala}



\begin{quote}
Before the new one: say the Arabic for drizzle, then the Arabic for a downpour.
\end{quote}

\subsection*{You'll want to know: \ar{هطل}}
\textbf{\ar{هطل}} — \emph{haṭala} — \textquotedblleft{}to pour down (of rain)\textquotedblright{}.

\begin{quote}
\ar{هطل} \ar{المطر} \ar{في} \ar{الليل}.
\end{quote}

\begin{quote}
\emph{haṭala al-maṭar fī al-layl.}
\end{quote}

\begin{quote}
The rain poured down in the night.
\end{quote}

Present: \ar{يهطل}, \emph{yahṭil}.

\begin{cousinweb}
\textbf{\ar{هطل}} (\emph{haṭala}) is built on the native Arabic root \emph{h-ṭ-l}, \textquotedblleft{}to pour (of water)\textquotedblright{} (Lane). It is the plain verb.
\end{cousinweb}

\begin{grammarlens}[title={What you've built}]
One more word for getting around.
\end{grammarlens}

\subsection*{Your turn}
\begin{itemize}
\raggedright
  \item \emph{Say it:} \emph{haṭala}
  \item \emph{Say it:} \emph{haṭala}, once more
  \item \emph{Say it:} \emph{haṭala al-maṭar fī al-layl.} — The rain poured down in the night.
  \item \emph{Recall:} say the Arabic for drizzle, then the Arabic for a downpour, then say \emph{haṭala} again
\end{itemize}

\begin{tcolorbox}[breakable,colback=teal!4,colframe=teal!35!black,title={Before you move on}]
What does \ar{هطل} mean? (\textquotedblleft{}To pour down (of rain)\textquotedblright{}.) Say it once more.
\end{tcolorbox}
\section[\texorpdfstring{\ar{صفا} (ṣafā)}{صفا (ṣafā)}]{\ar{صفا} (ṣafā) — to clear up; to become clear}

\label{lesson:AR-C481-safa}



\begin{quote}
Before the new one: say the Arabic for a downpour, then the Arabic for to pour down.
\end{quote}

\subsection*{You'll want to know: \ar{صفا}}
\textbf{\ar{صفا}} — \emph{ṣafā} — \textquotedblleft{}to clear up; to become clear\textquotedblright{}.

\begin{quote}
\ar{صفت} \ar{السماء} \ar{بعد} \ar{المطر}.
\end{quote}

\begin{quote}
\emph{ṣafat as-samāʾ baʿda al-maṭar.}
\end{quote}

\begin{quote}
The sky cleared after the rain.
\end{quote}

Present: \ar{يصفو}, \emph{yaṣfū}.

\begin{cousinweb}
\textbf{\ar{صفا}} (\emph{ṣafā}) is built on the native Arabic root \emph{ṣ-f-w}, \textquotedblleft{}to be clear, pure\textquotedblright{} (Lane). It is the plain verb.
\end{cousinweb}

\begin{grammarlens}[title={What you've built}]
One more word for getting around.
\end{grammarlens}

\subsection*{Your turn}
\begin{itemize}
\raggedright
  \item \emph{Say it:} \emph{ṣafā}
  \item \emph{Say it:} \emph{ṣafā}, once more
  \item \emph{Say it:} \emph{ṣafat as-samāʾ baʿda al-maṭar.} — The sky cleared after the rain.
  \item \emph{Recall:} say the Arabic for a downpour, then the Arabic for to pour down, then say \emph{ṣafā} again
\end{itemize}

\begin{tcolorbox}[breakable,colback=teal!4,colframe=teal!35!black,title={Before you move on}]
What does \ar{صفا} mean? (\textquotedblleft{}To clear up; to become clear\textquotedblright{}.) Say it once more.
\end{tcolorbox}
